% =========================================================
% Exemples File
% =========================================================

% Seções --------------------------------------------------

\section{Seção de Exemplo} % <-- Título da seção
\label{sec:e1}

    Uma seção se parece com isso.

\subsection{Subseção de Exemplo} % <-- Título da subseção
\label{sec:e2}

    Uma subseção se parece com isso.

\subsubsection{Subsubseção de Exemplo} % <-- Título da subsubseção
\label{sec:e3}

    Uma subsubseção se parece com isso.

\addcontentsline{toc}{section}{Seção Não Enumerada} % <-- Título da seção
\section*{Seção Não Enumerada} % <-- Título da seção

    Assim se parece uma seção não enumerada.
    Subseções e subsubseções também podem ser não enumeradas basta incluir o asterisco (*) no comando.
    As labels não funcionam bem em seções não enumeradas.

\section{Imagens}
\label{sec:imagens}

    A figura \ref{fig:e4} é um exemplo de imagem sozinha.
    A imagem pode ser referenciada com o comando \lstinline[language=TeX, style=TeX]|\ref{}|.

    \begin{figure}[htb]
        \centering
        \includegraphics[width=120pt]{images/figure.png}
        \caption{Legenda.}
        \label{fig:e4}
    \end{figure}

    As images também podem ser utilizadas lado a lado:

    \begin{figure}
        \centering
        \begin{minipage}{0.5\textwidth}
            \centering
            \includegraphics[width=120pt]{images/figure.png}
            \caption{Legenda figura PNG.}
            \label{fig:e6}
        \end{minipage}
        \begin{minipage}{0.5\textwidth}
            \centering
            %\includesvg[width=120pt]{images/figure.svg}
            \caption{Legenda figura SVG.}
            \label{fig:e7}
        \end{minipage}
    \end{figure}

\section{Código Fonte}
\label{sec:codigo}

    Texto texto texto texto texto texto texto texto texto texto texto
    \lstinline[language=Python, style=Python]|def incmatrix(genl1,genl2)|

\begin{lstlisting}[language=Python, style=Python]
import numpy as np 

def incmatrix(genl1,genl2):
    m = len(genl1)
    n = len(genl2)
    M = None #to become the incidence matrix
    VT = np.zeros((n*m,1), int)  #dummy variable
    
    #compute the bitwise xor matrix
    M1 = bitxormatrix(genl1)
    M2 = np.triu(bitxormatrix(genl2),1) 

    for i in range(m-1):
        for j in range(i+1, m):
            [r,c] = np.where(M2 == M1[i,j])
            for k in range(len(r)):
                VT[(i)*n + r[k]] = 1;
                VT[(i)*n + c[k]] = 1;
                VT[(j)*n + r[k]] = 1;
                VT[(j)*n + c[k]] = 1;
                
                if M is None:
                    M = np.copy(VT)
                else:
                    M = np.concatenate((M, VT), 1)
                
                VT = np.zeros((n*m,1), int)
    
    return M
\end{lstlisting}
